\documentclass{article}
\usepackage{amsmath,amssymb}
\usepackage{xurl}
\usepackage[hidelinks]{hyperref}
% Shared commands for notes in docs/paper-gaps/.

\DeclareUnicodeCharacter{2080}{\ifmmode{}_{0}\else\textsubscript{0}\fi}
\DeclareUnicodeCharacter{2081}{\ifmmode{}_{1}\else\textsubscript{1}\fi}
\DeclareUnicodeCharacter{2082}{\ifmmode{}_{2}\else\textsubscript{2}\fi}
\DeclareUnicodeCharacter{2099}{\ifmmode{}_{n}\else\textsubscript{n}\fi}

\newcommand{\Lean}{\textsc{Lean}}
\ExplSyntaxOn
\NewDocumentCommand{\leanid}{m}{
  \ifmmode
    \textsf{\detokenize{#1}}
  \else
    \group_begin:
    \urlstyle{sf}
    \tl_set:Nn \l_tmpa_tl {#1}
    \tl_replace_all:Nnn \l_tmpa_tl {\_} {_}
    \exp_args:NV \path \l_tmpa_tl
    \group_end:
  \fi
}
\ExplSyntaxOff
\providecommand{\tr}{\operatorname{tr}}
\newcommand{\tnleanIssueBase}{https://github.com/LionSR/TNLean/issues/}
\newcommand{\tnleanPrBase}{https://github.com/LionSR/TNLean/pull/}
\newcommand{\tnleanDocBase}{https://sirui-lu.com/TNLean/docs/}
\newcommand{\ghissue}[1]{\href{\tnleanIssueBase#1}{GitHub issue \##1}}
\newcommand{\ghpr}[1]{\href{\tnleanPrBase#1}{PR \##1}}
\newcommand{\doclink}[1]{\href{\tnleanDocBase#1.html}{\path{#1}}}

% Dirac notation and the identity operator, as in blueprint/src/macros/common.tex.
% \providecommand keeps note-local definitions authoritative.
\usepackage{dsfont}
\providecommand{\ket}[1]{|#1\rangle}
\providecommand{\bra}[1]{\langle #1|}
\providecommand{\braket}[2]{\langle #1 | #2 \rangle}
\providecommand{\ketbra}[2]{|#1\rangle\!\langle #2|}
\providecommand{\Id}{\mathds{1}}
\providecommand{\one}{\mathds{1}}
\providecommand{\C}{\mathbb{C}}
\providecommand{\MN}[1]{M_{#1}(\C)}
\providecommand{\MD}{M_D(\C)}

% Verdict marker: \gapnote{<kind>}{<status>}. Kind is the policy
% classification (clarification, local-correction, scope-restriction,
% unfaithful, false-source, open-gap); status is open, wip (elimination
% underway), resolved, or historical. Machine-read by the site index;
% prints a verdict line.
\newcommand{\gapnote}[2]{\par\noindent\textbf{Verdict:} #1 (#2).\par}

\title{Asymptotic Local Channel Conversion without Classical Feedforward}
\author{}
\date{2026-10-04}
\begin{document}
\maketitle
\gapnote{scope-restriction}{open}

\section{The source relation}
In the paragraph ``Phases of matter'' of~\cite{Piroli2021LOCC}, a sequence
$\{\ket{\psi_N}\}$ is convertible into $\{\ket{\varphi_N}\}$ if there is a
fixed number $k$ of compositions of channels assisted by measurements and
classical communication, each of depth $f(N)$, whose outputs $\sigma_N$
satisfy
\begin{align}
    \lVert\sigma_N-\ket{\varphi_N}\bra{\varphi_N}\rVert_1 & \longrightarrow 0.
\end{align}
Here $f$ is polylogarithmic. Mutual convertibility defines the source's
phase equivalence. The fixed number of compositions and the polylogarithmic
depth bound are allowed to depend on the pair of sequences, but not on $N$.

\section{The enlarged-site local-channel model}
The available local protocols alternate onsite channels with layers of
channels on disjoint nearest-neighbor pairs. They act on arbitrary input
operators and allow changes of local dimension. A layer of two-site channels
on the resulting enlarged sites costs one depth unit, regardless of that
dimension. Intermediate dimensions may depend on the system size without a
uniform bound. The model has no classical feedforward between separated
sites. For this class, approximate conversion
means $\lVert\Psi(\rho)-\sigma\rVert_1\le\varepsilon$, with protocol depth at
most $T$.\footnote{Formal declaration:
    \leanid{QuantumCircuit.IsApproxLocalChannelConversion} in
    \doclink{TNLean/Circuit/Channel/ApproximateConversion}.}
For Hermitian input operators, composition adds the depths and errors:
\begin{align}
    \lVert\Psi_2(\Psi_1(\rho))-\tau\rVert_1
     & \le \lVert\Psi_2(\Psi_1(\rho)-\sigma)\rVert_1
    +\lVert\Psi_2(\sigma)-\tau\rVert_1               \\
     & \le \lVert\Psi_1(\rho)-\sigma\rVert_1
    +\lVert\Psi_2(\sigma)-\tau\rVert_1
    \le\varepsilon_1+\varepsilon_2.
\end{align}
The middle inequality is trace-norm contractivity for trace-preserving
positive maps on Hermitian operators.

For families, the error tends to zero and the depth is eventually bounded
by $C(1+\log(N+1))^p$, for a nonnegative constant $C$ and integer $p\ge0$.
Mutual conversion is an equivalence relation on eventually Hermitian
families, and in particular on families of density matrices.\footnote{Formal
    declarations: \leanid{QuantumCircuit.IsAsymptoticLocalChannelConversion} and
    \leanid{QuantumCircuit.IsLocalChannelPhaseEquivalent.equivalence} in
    \doclink{TNLean/Circuit/Channel/AsymptoticConversion}.}
The shift in the logarithm only affects the bound at small system sizes.
Finite compositions are already included: their depths add, and a finite sum
of polylogarithmic bounds is polylogarithmic. No bound on the dimension of
onsite ancillas is imposed.

This gate convention differs from the source. In the Supplement's proof of
Proposition~2 (p.~8), the intersite gates act only on the physical qudits;
the entanglement bound across a cut depends on their fixed dimension $d$,
not on the number of onsite ancillas. Allowing one intersite layer to act
on an arbitrarily enlarged local space does not by itself give a
uniform-depth simulation in that convention. No such simulation or
inclusion between the two asymptotic conversion relations is proved here.

\section{Blocking sites}
In the proof of its classification theorem for matrix product states, the
source groups $q$ neighbouring sites of a chain of $N$ sites into one site of
dimension $d^q$, giving a chain of $N/q$ sites~\cite[Supplemental Material;
    \texttt{main.tex} of arXiv:2103.13367v3, line 287]{Piroli2021LOCC}. An
operation on the blocked chain acts on $q$ adjacent qudits of the original
chain, and each one- or two-site gate on the blocked chain is realized by
fewer than $2q^2$ sequential two-site gates on the original
chain~\cite[\texttt{main.tex}, line 289]{Piroli2021LOCC}.

Blocking itself is a relabelling of the computational basis. It is a
unitary channel and is formalized as a step of depth zero. A blocked
conversion first blocks $k$ sites and then runs a local channel protocol of
depth at most $T$ on the chain of blocks.\footnote{Formal declaration:
    \leanid{QuantumCircuit.IsBlockedLocalChannelConversion} in
    \doclink{TNLean/Circuit/Channel/SiteBlocking}.}
The blocking factor $k$ is part of the relation: blocking the whole chain
into one site and applying a replacement channel there would convert every
density matrix into every other one at depth zero. As in the source, the
blocking factor is positive. Locality is transported by
rescaling: the dual of such a conversion maps an operator acting on a set
$X$ of blocks to an operator acting within distance $kT$ of the sites of
$X$ in the original chain.

As in the enlarged-site model above, one layer of two-site channels on the
blocked chain costs one depth unit. The source's compilation of a blocked
gate into fewer than $2q^2$ nearest-neighbour gates on the original chain is
not formalized, so no bound on the original-chain depth of a blocked
conversion is asserted.

\section{Regrouping an existing local channel protocol}
For $M,k,d\ge1$ there is an internal comparison in the enlarged-site model: a depth-$T$
protocol on $Mk$ sites can be regrouped into $M$ sites with exactly the
conjugated channel on every input operator. The identity also holds with an
arbitrary untouched external reference. If every intermediate local dimension
is positive and at most $B$, every regrouped dimension is positive and at most
$B^k$.\footnote{Formal
    declarations: \leanid{QuantumCircuit.IsDimensionBoundedLocalChannelProtocol.exists_regrouped}
    and \leanid{QuantumCircuit.regroupChannel_blockSites_reference} in
    \doclink{TNLean/Circuit/Channel/RegroupedProtocol}.}

The resulting depth is at most $2T$ for even $M$ and $3T$ in general.
The original gates whose left sites belong to one block are composed into
one channel on that block and its successor. Alternating bond colors suffice
on an even ring; an odd ring uses a third color for its closing bond.
For $M=3,k=2$, the disjoint original bonds $(1,2),(3,4),(5,0)$ become a
triangle. Thus the direct partition into whole-block matching layers need
not use only two layers. This does not exclude other implementations of the
same channel. Onsite operations remain free; when $M=1$ the entire regrouped
channel has depth zero. The $M=2$ bound is $2T$.

Consequently, an unblocked conversion of depth at most $T$ followed by a
blocked conversion of depth at most $U$ composes with depth bound $2T+U$
for even $M$, or $3T+U$ in general.\footnote{Formal declaration:
    \leanid{QuantumCircuit.IsLocalChannelConversion.then_blocked}.}
This comparison implements grouping of an existing protocol. It supplies
neither the reverse fixed-physical-port compiler nor the source's fixed
number of classical-communication blocks, and does not prove the phase
classification.

\section{Remaining scope}
The equivalence proved here belongs to the enlarged-site local-channel model.
It is not identified with the source's phase equivalence or asserted to be
a subclass at the same asymptotic depth. The classification asserting that
every translation-invariant matrix product state belongs to the trivial
source phase is not obtained. Finite adaptive channels on arbitrary input
operators and their additive-depth composition are now available separately,
as described in~\cite{gap:psc21_adaptive_channel_round_scope}. They are not
included in the non-feedforward asymptotic relation formalized here.

The remaining source comparison needs a decomposition into a fixed number
of \textit{QCcc} blocks, with that number independent of system size. Each
block must satisfy the source's internal measurement/control order and
one-measurement-per-site rule. The finite adaptive tree does not record this
decomposition: counting its individual instrument steps would not express
the source condition. The comparison must also distinguish fixed-dimensional
physical qudits from onsite ancillas and control the intersite gate cost
uniformly in the chain length. A proved simulation bound or explicit
restrictions on intermediate dimensions could supply this resource comparison.
Once those source-model depth bounds are established, trace-norm contractivity
supplies the same error estimate. The separate restrictions of the
matrix-product-state classification are recorded
in~\cite{gap:psc21_mps_classification_scope}.

\bibliographystyle{alpha}
\bibliography{references}
\end{document}
